\section{Introducción}
Esta clase comienza con Xinyun Chen mencionando que «\textit{2025 will be the year of Agents}», y explicando que este semestre el curso Advanced LLM Agents SP25 de Berkeley RDI, en colaboración con Google DeepMind y Cal, cuenta tanto con casos consolidados en el track de aplicaciones como con un nuevo track de investigación centrado en inference-time reasoning. El docente enfatiza que más de 15,000 inscritos al MOOC, más de 3,000 participantes en el hackathon y más de 200,000 dólares en recursos de la industria buscan mostrar que este curso no es una demostración puramente teórica, sino que de verdad pretende implementar las técnicas de tiempo de inferencia en escenarios de ingeniería complejos.

\begin{knowledgebox}{Posicionamiento del curso}
De manera modular, se divide la técnica de tiempo de inferencia en tres partes: «prompt + búsqueda + iteración». El docente compartirá en esta clase, con detalle, la estrategia, las plantillas de instrucciones y los puntos clave de ingeniería de cada una de ellas, de modo que, bajo el tema de «invertir cómputo en el tiempo de inferencia», atendamos a la vez las tres dimensiones de tokens, ancho y profundidad.
\end{knowledgebox}

\subsection{Resumen de la sección}
La introducción establece el contexto de arranque en frío de los Agentes y enfatiza el inference-time compute como una palanca que mejora el razonamiento sin modificar los pesos, organizando el contenido en tres partes: prompt, búsqueda e iteración; las secciones siguientes se desarrollan conforme a esta lógica.
